\documentclass[12pt,a4paper]{article}
\usepackage[utf8]{inputenc}
\usepackage[vietnamese]{babel}
\usepackage[left=2cm,right=2cm,top=2cm,bottom=2cm]{geometry}
\usepackage{mathptmx}
\usepackage{enumitem}

\pagestyle{plain}
\linespread{1.15}

\begin{document}

% --- HEADER QUỐC HIỆU & CƠ QUAN BAN HÀNH ---
\noindent
\begin{minipage}[t]{0.45\textwidth}
    \begin{center}
        \small
        CÔNG AN THÀNH PHỐ HÀ NỘI\\
        \textbf{\underline{PHÒNG CẢNH SÁT HÌNH SỰ}}\\[0.1cm]
        Số: 14/BB-KXNƠ\\
        \textit{V/v khám xét khẩn cấp chỗ ở và thân thể}
    \end{center}
\end{minipage}
\hfill
\begin{minipage}[t]{0.52\textwidth}
    \begin{center}
        \small
        \textbf{CỘNG HÒA XÃ HỘI CHỦ NGHĨA VIỆT NAM}\\
        \textbf{\underline{Độc lập - Tự do - Hạnh phúc}}\\[0.1cm]
        \textit{Hà Nội, ngày 25 tháng 07 năm 2026}
    \end{center}
\end{minipage}

\vspace{0.6cm}

% --- TIÊU ĐỀ VĂN BẢN ---
\begin{center}
    \textbf{\large BIÊN BẢN KHÁM XÉT KHẨN CẤP CHỖ Ở VÀ THÂN THỂ}\\
    \textit{(Đối tượng Trần Thị Hà — Số 8, Ngõ 12 Đường Bờ Sông)}
\end{center}

\vspace{0.3cm}

\noindent Căn cứ Lệnh khám xét khẩn cấp số 22/LKX-CSHS ngày 25/07/2026;\\
Vào hồi 13 giờ 00 phút ngày 25 tháng 07 năm 2026, tại phòng trọ số 8, Ngõ 12 Đường Bờ Sông, Phân khu Cảng.\\
Tổ công tác Phòng CSHS chủ trì cùng sự chứng kiến của Chủ nhà trọ và Cán bộ phụ nữ phường tiến hành khám xét nơi ở và thân thể đối tượng **TRẦN THỊ HÀ (Sinh năm 1994)**. Kết quả thu giữ:

\section*{I. ĐỒ VẬT, TÀI LIỆU THU GIỮ TẠI PHÒNG TRỌ}
\begin{enumerate}[label=\arabic*., leftmargin=1.5cm]
    \item \textbf{01 Chiếc áo khoác gió mỏng màu xám đen có mũ trùm đầu:} Treo sau cánh cửa phòng trọ, form dáng nhỏ gọn (size S), nồng nặc mùi tinh dầu hoa nhài. Dưới gấu áo dính bụi đất mùn đặc trưng quanh gốc cây xoan trước ngõ nhà Khang.
    \item \textbf{01 Cuộn len màu đỏ thẫm đan dở:} Cùng que đan bằng inox đặt trên bàn làm việc giữa thời tiết nắng nóng 39°C.
    \item \textbf{01 Vỉ thuốc an thần Sedative nhãn hiệu Diazepam 5mg:} Đã bóc dở 4 viên, trùng khớp với hoạt chất an thần tìm thấy trong cặn ấm trà hoa cúc và dịch dạ dày nạn nhân.
    \item \textbf{01 Cuốn sổ nhật ký bìa hồng:} Ghi chép chi tiết lịch trình sinh hoạt từng giờ, số đo kích thước cúc áo, độ dài tóc mai và các mối quan hệ xã hội của nạn nhân Nguyễn Văn Khang.
\end{enumerate}

\section*{II. VẬT CHỨNG ĐẶC BIỆT THU GIỮ TRÊN THÂN THỂ ĐỐI TƯỢNG}
Khi cán bộ nữ tiến hành khám xét thân thể đối tượng Trần Thị Hà, phát hiện:
\begin{itemize}[leftmargin=1.5cm]
    \item \textbf{Dấu vết thương tích trên bàn tay phải:} Lòng bàn tay và kẽ ngón trỏ phải của Hà có 01 vết rách da sâu dài 2.5cm, rỉ máu rớm miệng vết thương, đã được dán tạm 02 miếng băng dán y tế cá nhân Urgo (hoàn toàn phù hợp với cơ chế bị cạnh sắc mảnh vỡ thủy tinh cứa ngược khi dùng lực đâm).
    \item \textbf{Vật chứng giấu trong khe áo ngực:} Thu giữ 01 gói giấy ăn gấp gọn, bên trong bọc **01 lọn tóc mai dính máu khô** được cắt bằng kéo bấm chỉ.
    \item \textbf{Kết quả giám định sinh học khẩn cấp (Viện Khoa học Hình sự):} Dấu vết máu trên lọn tóc mai trùng khớp 100\% mẫu ADN của nạn nhân Nguyễn Văn Khang.
\end{itemize}

\vspace{0.8cm}

% --- CHỮ KÝ NGHỊ ĐỊNH 30 ---
\noindent
\begin{minipage}[t]{0.45\textwidth}
    \small
    \textbf{\textit{Nơi nhận:}}\\
    - Viện KSMND TP. Hà Nội;\\
    - Hồ sơ chuyên án TEST-99;\\
    - Lưu: VT, CSHS.
\end{minipage}
\hfill
\begin{minipage}[t]{0.5\textwidth}
    \begin{center}
        \small
        \textbf{CÁN BỘ CHỦ TRÌ KHÁM XÉT}\\
        \textit{(Ký, ghi rõ họ tên)}\\[1.5cm]
        \textbf{Đại úy Lê Minh}
    \end{center}
\end{minipage}

\end{document}
